%Intro
OpenCV\cite{opencv} (Open Source Computer Vision) is a library that implements state of the art real-time computer vision 
algorithms. 

It is cross-platform and it is released under a BSD\cite{BSD} license. Previously it was integrated in ROS \cite{ros}, but
now it is used as a stand-alone library. This allows to easily use within ROS the latest OpenCV version. 

\begin{figure}[h]
	\begin{center}
    \includegraphics[scale=1]{img/opencv/logo.eps}
	\caption[OpenCV Logo]{OpenCV Logo}
	\end{center}
\end{figure}

Those qualities described above as well as its optimization of various algorithms make this library an extremely useful 
asset for robotic vision projects. 

% In this project's context

In this project the C++ interface of the OpenCV library has been used. The functionalities provided by this library 
are used in all the 2D information management: 
\begin{itemize}
	\item{Segmentation of the ROI (Region Of Interest): } This library allows to crop the original raw image coming
	from the RGB-D sensor. In this application, the region cropped is the one around the hand being used. 
	\item {Keypoints and descriptors extraction: } The OpenCV library implements multiple optimized state of the art
	algorithms. Among them, there are a number of descriptor extractors such as SIFT, SURF, ORB. This last one is 
	being used in this project due to its open-sourceness, and robustness/speed relation. 
	\item {Descriptors matching: } There are many different matchers that might be used. In this project, a FlannBasedMatcher
	is being used for this purpose. 
\end{itemize}

More information about the reasons of chosing those specific classes and algorithms from the OpenCV library might be found
in next chapters. 

